\begin{explanationbox}
	\textbf{\textcolor{CExplanation}{一句话直觉}}
	将斜边上的正方形分割成四个直角三角形和一个小正方形，通过面积相等列出方程，直接得出勾股定理。

	\medskip
	\textbf{\textcolor{CExplanation}{核心拆解}}
	\begin{description}[style=nextline, leftmargin=2.8em, labelsep=0.5em, itemsep=0.45em]
		\item[\textbf{要点一（图形构成）：}]
			直角三角形的直角边为 $a,b$（$a\ge b$），斜边 $c$。四个全等直角三角形按旋转对称摆放后，它们的外轮廓构成一个边长为 $c$ 的大正方形，内部留下一个边长为 $a-b$ 的小正方形。
		\item[\textbf{要点二（面积关系）：}]
			每个直角三角形面积为 $\frac12 ab$，四个总面积为 $2ab$；小正方形面积为 $(a-b)^2$。大正方形面积为 $c^2$。
	\end{description}

	\medskip
	\textbf{\textcolor{CExplanation}{几何本质（出入相补）}}
	平面图形的面积在切割平移下保持不变。弦图将斜边正方形分割后，各部分可以重新拼成两个分别以 $a$ 和 $b$ 为边长的正方形，从而揭示三边关系。

	\medskip
	\textbf{\textcolor{CExplanation}{代数意义（恒等变形）}}
	由面积等式 $c^2 = 2ab + (a-b)^2$，展开右边得 $2ab + a^2 -2ab + b^2 = a^2 + b^2$，所以 $c^2 = a^2+b^2$。代数上就是完全平方公式的逆向使用。

	\medskip
	\textbf{\textcolor{CExplanation}{考点价值（考试方向）}}
	\begin{itemize}[leftmargin=*, itemsep=0.5em, topsep=0.4em]
		\item \textbf{考法一（勾股计算）：} 直接给出直角三角形的两边长，求第三边；或通过弦图条件列方程求边长。
		\item \textbf{考法二（面积思想）：} 利用弦图中的面积关系推导代数恒等式，训练数形结合能力。
	\end{itemize}

	\medskip
	\textbf{\textcolor{CExplanation}{顿悟点}}
	\textbf{面积是连接几何与代数的桥梁。} 观察弦图，我们发现同一个正方形的面积有两种表达方式，从而将未知量 $c$ 与已知量 $a,b$ 联系起来。这种“面积两次表达”的方法是方程思想的典型应用。

	\medskip
	\textbf{\textcolor{CExplanation}{使用场景}}
	\begin{itemize}[leftmargin=*, itemsep=0.5em, topsep=0.4em]
		\item \textbf{场景一（定理证明）：} 在没有数具的情况下，仅靠几何构造即可证明勾股定理，适合课堂引入。
		\item \textbf{场景二（数形结合）：} 沟通完全平方公式 $(a\pm b)^2 = a^2 \pm 2ab + b^2$ 与勾股定理，加深对代数与几何联系的理解。
	\end{itemize}

\end{explanationbox}
